\begin{abstract}
机器学习系统（ML systems）方向的面试，不论是训练 infra、推理引擎、算子与编译器还是 AI 芯片岗，考的几乎是同一种能力：
给定模型、硬件和负载，能不能在几分钟内用第一性原理估出时间、显存和通信的下界，指出瓶颈在哪里，并说清楚换一种设计会改变什么。
本文把 Stanford CS149（并行计算）、Stanford CS229S（机器学习系统）和 CMU 15-849（机器学习系统研讨课）三门课的内容，
围绕「算得过来吗、放得下吗、传得过来吗」这三个问题重新组织成 18 个模块。内容从性能模型、并行硬件与 GPU 编程模型讲起，
经过 GEMM 与 FlashAttention 等算子、专用加速器、ML 编译器与框架，一直到分布式训练（数据、张量、流水线、序列并行与 ZeRO）、
LLM 推理与服务（KV cache、PagedAttention、chunked prefill、PD 分离、推测解码、量化）、模型效率（LoRA、MoE、状态空间模型、
scaling law）以及集群调度与可靠性。和一般的综述不同，本文的每个结论都配有能复算的白板估算。数字只有三个来源：正文里的推导，
配套计算器 \texttt{tools/mlsys\_calc.py} 的输出，以及原论文和本文配套的 9 个 CPU 实验的原始 CSV；文中始终区分「模型输出」和「GPU 实测」。
全文有 57 个「面试追问」框（每个都附常见错答）、16 个白板估算、15 幅图和 19 张表。此外，本文给出一个把工程能力和研究能力分开刻画的能力矩阵、
一个「结论、估算、瓶颈、权衡、证据」五步答题模板，以及把课程实验变成面试故事和研究练习的方法。文中也纠正了几个流传很广的误解：
FlashAttention 的 HBM 流量在序列长度上仍是 $\Theta(T^2 d^2/M)$，并没有降到线性；1F1B 并不减少流水线气泡；
ZeRO-3 的通信量是数据并行的 1.5 倍，而不是 3 倍。
\end{abstract}

\begin{IEEEkeywords}
机器学习系统，roofline，FlashAttention，分布式训练，ZeRO，张量并行，流水线并行，LLM 推理，KV cache，
PagedAttention，推测解码，量化，ML 编译器，面试，能力建设
\end{IEEEkeywords}
